\section{独立任务分区与资源配置}
\subsection{固定联合计划下的分区模型}
Q4 保持 Q3 的组批、访问顺序、开始时刻与通信保障关系不变。把出现在同一运输架次上的服务区连接，得到四个不可拆分服务区连通分量。任务组只能由这些分量合并，否则会把一个既定运输架次拆给两个组。

对 $K=2,3$，分别穷举四个分量的无标号集合划分，候选数为 7 和 6。设组 $k$ 内某类资源占用区间集合为 $\mathcal I_{kh}$，则在允许同类资源重新编号时所需数量为
\begin{equation}
n_{kh}=\max_t\sum_{I\in\mathcal I_{kh}}\mathbf 1\{t\in I\},\qquad
\Delta_h=\max\left(0,\sum_kn_{kh}-I_h^{\rm stock}\right).
\label{eq:partition}
\end{equation}
采用半开区间处理任务衔接，无人机、电池、中继及能源组件的占用终点分别纳入返航、充电和周转规则。相同时间容差用于求解与验证，避免把浮点误差计为额外资源。

若一个中继架次为不同任务组的运输架次提供服务，为保持各组独立执行，我们在各组复制该中继任务的完整位置与服务安排，并分别计数。这个解释保留原来的通信关系但增加设备需求；若要求单个中继任务也不可跨组复制，则须进一步合并相应服务区，属于另一种分区约束，本文未将两者混为同一结果。

划分按“库存缺口总数、资源需求总数、组间工作量差异”的词典序选择。不同类别的数量直接相加是配置规模指标，不是经济成本；没有设备单价数据时不称为成本最优。工作量由运输占用时间加中继通信服务时间构成，既不同于完工历时，也不是人员工时。

\subsection{资源需求、冗余与决策含义}
\textbf{固定联合方案独立划为两组时总资源缺口为 \GapTwo{}，三组时为 \GapThree{}。}表\ref{tab:resources}给出每类库存与需求，分区对应关系见图\ref{fig:partition}。分组使原本可跨任务共享的闲置资源无法互补，因此总体任务量不变也可能增加峰值配置。
\paperfig{fig09_partition}{独立分区增加固定方案的资源需求。灰色竖线为现有库存，圆点和方点分别为两组、三组方案；下图显示服务区归属，纹理与组号提供灰度可辨识编码。缺口总数按台或组件计数，不代表按价格加权的成本。}{partition}
\papertable{resource_gaps}{固定 Q3 计划独立执行时的库存、需求与缺口}{resources}

两组方案保留了更多组内复用机会。三组方案隔离程度更高，但更多的资源需求不能由现有库存覆盖。资源配置选择因而不能只看组数或每组服务区数量，还要看具体时间上的重叠与能源周转。这里的最优分区是相对于固定 Q3 方案和给定比较准则的精确枚举结果，不是允许重新调度情况下的最小资源需求。
\FloatBarrier
